%==============================================================================
% PARTIDA 03: CANALIZACIONES, TUBERIAS Y CONDUCTOS
% Codigos: OE.5.2.2.1 al OE.5.2.2.7
%==============================================================================
\subsection{OE.5.2.2.1 Tuberia PVC Liviana (PV-L) 3/4'' -- Canalizacion de Alumbrado}

\subsubsection{Especificaciones Tecnicas}

Tuberia de policloruro de vinilo (PVC) tipo liviano para canalizacion electrica de los circuitos de alumbrado C1, C3 y C5. Debe cumplir con la NTP 399.006 (Tubos de PVC para instalaciones electricas) y la NTP-IEC 61386 (Sistemas de tubos para la gestion del cableado).

\begin{itemize}[nosep]
  \item \textbf{Tipo:} PVC liviano (PV-L), clase de presion liviana.
  \item \textbf{Diametro nominal:} 3/4'' (20 mm de diametro exterior nominal).
  \item \textbf{Diametro interior minimo:} 16.5 mm.
  \item \textbf{Espesor de pared nominal:} 1.5 mm ± 0.2 mm (segun NTP 399.006 Tabla 1).
  \item \textbf{Color:} Naranja o crema (estandar para canalizaciones electricas).
  \item \textbf{Longitud util:} 3.00 m por barra (tolerancia: +0, -10 mm).
  \item \textbf{Temperatura de servicio:} -5 °C a +60 °C.
  \item \textbf{Resistencia al impacto:} 2.0 J a 23 °C (segun NTP 399.006).
  \item \textbf{Resistencia a la compresion radial:} 250 N minimo.
  \item \textbf{Grado de inflamabilidad:} Autoextinguible (V-0 segun UL 94).
\end{itemize}

\paragraph{Condiciones de instalacion:}
\begin{itemize}[nosep]
  \item Maximo 3 curvas de 90° entre cajas de derivacion consecutivas (CNE-U Articulo 350.12).
  \item Distancia maxima entre cajas de paso o derivacion: 15 m (CNE-U Articulo 240.4).
  \item Porcentaje maximo de ocupacion para conductores: 40\% del area transversal (CNE-U Tabla 350.1).
  \item Para 3 conductores de 1.5 mm2 (diametro exterior ≈ 3.5 mm), area ocupada: 3 x 9.6 mm2 = 28.8 mm2. Area interior de tuberia 3/4'': 213.8 mm2. Porcentaje de ocupacion: 13.5\% < 40\% OK.
\end{itemize}

\subsubsection{Requisitos de Materiales}

\begin{longtable}{L{5.0cm} L{2.0cm} L{7.5cm}}
\toprule
\textbf{Material} & \textbf{Cant.} & \textbf{Especificacion tecnica} \\
\midrule
Tuberia PVC liviana (PV-L) 3/4'' x 3 m & 40 barras & NTP 399.006, espesor 1.5 mm, color naranja \\
Curva PVC liviana 3/4'' 90° & 12 pza | Radio medio 50 mm, para cambios de direccion \\
Curva PVC liviana 3/4'' 45° & 3 pza | Radio medio 50 mm, para cambios suaves \\
Union universal PVC liviana 3/4'' & 20 pza | Para empalme de barras de tuberia \\
Pegamento PVC (cemento solvente) & 1/4 gln | Para PVC rigido, secado rapido (30 s), color transparente \\
Abrazadera tipo gancho 3/4'' & 40 pza | Para fijacion a losa antes del vaciado \\
Guia de alambre galvanizado calibre 14 & 50 m | Para paso de conductores despues de instalada la canalizacion \\
\bottomrule
\end{longtable}

\subsubsection{Metodo de Ejecucion}

\begin{enumerate}
  \item \textbf{Trazo de rutas:} Marcar en la losa de techo el recorrido de las canalizaciones de alumbrado segun planos IE-02, IE-03, IE-04. Las rutas deben ser ortogonales (paralelas a muros), minimizando cruces y cambios de direccion.
  \item \textbf{Corte de tuberia:} Medir y cortar la tuberia PVC a las longitudes requeridas utilizando segueta de diente fino (32 dientes/pulgada). El corte debe ser perpendicular al eje del tubo. Eliminar rebabas con lima.
  \item \textbf{Ensamblaje:}
  \begin{enumerate}
    \item Limpiar el extremo del tubo y el interior del accesorio con trapo humedo en thinner.
    \item Aplicar pegamento PVC en el extremo del tubo (2/3 de la longitud del accesorio) y en el interior del accesorio.
    \item Insertar el tubo en el accesorio hasta el tope, girando 1/4 de vuelta para distribuir el pegamento.
    \item Mantener la presion durante 15 segundos (tiempo de fraguado inicial).
    \item Tiempo de curado completo: 4 horas antes de someter a carga.
  \end{enumerate}
  \item \textbf{Fijacion a losa:}
  \begin{enumerate}
    \item Colocar la tuberia sobre la armadura de la losa.
    \item Fijar con alambre calibre 16 amarrado a la malla de acero, o con abrazaderas tipo gancho cada 0.80 m como maximo.
    \item La separacion entre la tuberia y el encofrado debe ser tal que quede embebida dentro del concreto (recubrimiento minimo: 15 mm).
  \end{enumerate}
  \item \textbf{Prueba de paso:} Antes del vaciado de concreto, pasar una guia de alambre calibre 14 a traves de cada tramo de tuberia para verificar que no existan obstrucciones, aplastamientos o pegamento obstruyendo la seccion interior.
  \item \textbf{Proteccion durante vaciado:} Sellar los extremos de la tuberia con cinta aislante para evitar el ingreso de lechada de concreto.
\end{enumerate}

\subsubsection{Control de Calidad}

\begin{longtable}{L{4.5cm} L{4.5cm} L{4.5cm}}
\toprule
\textbf{Parametro} & \textbf{Metodo de verificacion} & \textbf{Criterio de aceptacion} \\
\midrule
Diametro interior & Calibrador Vernier & Minimo 16.5 mm \\
Espesor de pared & Calibrador Vernier & 1.5 mm ± 0.2 mm \\
Pendiente & Nivel de burbuja en tramos horizontales & Minimo 0.5\% \\
Separacion entre soportes & Cinta metrica & Maximo 0.80 m \\
Continuidad de canalizacion & Prueba de paso con guia calibre 14 & Paso libre sin obstrucciones \\
Curvas maximo permitidas & Inspeccion visual del recorrido & Maximo 3 curvas de 90° entre cajas \\
Porcentaje de ocupacion & Calculo segun CNE-U Tabla 350.1 & Maximo 40\% del area transversal \\
\bottomrule
\end{longtable}

\subsubsection{Metodo de Medicion}

Se medira por metro lineal (m) de tuberia PVC liviana de 3/4'' instalada, incluyendo curvas, uniones, accesorios y fijacion. No se descuentan curvas ni accesorios; la longitud se mide sobre el eje de la canalizacion.

\subsubsection{Unidad de Medida}

m (metro lineal).

\subsubsection{Forma de Pago}

Pago por metro lineal instalado y verificado mediante prueba de paso antes del vaciado de concreto. Aprobado por supervision.

%------------------------------------------------------------------------------
\subsection{OE.5.2.2.2 Tuberia PVC Pesada (PV-P) 3/4'' -- Canalizacion de Tomacorrientes}

\subsubsection{Especificaciones Tecnicas}

Tuberia de PVC tipo pesado para canalizacion de circuitos de tomacorrientes C2, C4 y C6. Debe cumplir con NTP 399.006.

\begin{itemize}[nosep]
  \item \textbf{Tipo:} PVC pesado (PV-P), clase de presion pesada.
  \item \textbf{Diametro nominal:} 3/4'' (20 mm exterior).
  \item \textbf{Diametro interior minimo:} 15.0 mm.
  \item \textbf{Espesor de pared nominal:} 2.0 mm ± 0.2 mm.
  \item \textbf{Color:} Naranja.
  \item \textbf{Resistencia al impacto:} 4.0 J a 23 °C (mayor resistencia que PV-L).
\end{itemize}

\paragraph{Ocupacion de conductores en PV-P 3/4'':}
Para 3 conductores de 2.5 mm2 (diametro exterior ≈ 4.2 mm): area ocupada por conductor = 13.9 mm2. Area total 3 conductores = 41.6 mm2. Area interior de tuberia 3/4'' PV-P = 176.7 mm2. Porcentaje = 23.5\% < 40\% OK.

El metodo de ejecucion, control de calidad y forma de pago son identicos a la partida OE.5.2.2.1.

\subsubsection{Metodo de Medicion}

Se medira por metro lineal (m) de tuberia PVC pesada de 3/4'' instalada.

\subsubsection{Unidad de Medida}

m (metro lineal).

%------------------------------------------------------------------------------
\subsection{OE.5.2.2.3 Tuberia PVC Pesada (PV-P) 1'' -- Canalizacion Alimentador General}

\subsubsection{Especificaciones Tecnicas}

Tuberia PVC pesada de 1'' (25 mm diametro exterior) para el alimentador general que transporta conductores de 10 mm2 desde la caja portamedidor hasta TG-01.

\begin{itemize}[nosep]
  \item \textbf{Diametro exterior nominal:} 25 mm.
  \item \textbf{Diametro interior minimo:} 20 mm.
  \item \textbf{Espesor de pared:} 2.5 mm ± 0.2 mm.
  \item \textbf{Ocupacion:} Para 3 conductores de 10 mm2 (diametro ≈ 6.5 mm): area ocupada total = 99.5 mm2. Area interior = 314.2 mm2. Porcentaje = 31.7\% < 40\% OK.
\end{itemize}

\paragraph{Metodo de ejecucion:} Identico a OE.5.2.2.1, con pendiente minima del 0.5\% hacia la caja portamedidor.

\subsubsection{Metodo de Medicion}

Se medira por metro lineal (m) de tuberia PVC pesada de 1'' instalada.

\subsubsection{Unidad de Medida}

m (metro lineal).

%------------------------------------------------------------------------------
\subsection{OE.5.2.2.4 a OE.5.2.2.7 -- Curvas y Accesorios de Canalizacion}

\subsubsection{Especificaciones Tecnicas}

Las curvas PVC permiten cambios de direccion en las canalizaciones manteniendo la seccion interior libre. Deben ser del mismo tipo y diametro que la tuberia correspondiente. Los accesorios de canalizacion comprenden el pegamento, conectores, uniones y demas insumos necesarios para completar la instalacion.

\begin{longtable}{L{5.0cm} L{3.0cm} L{6.5cm}}
\toprule
\textbf{Codigo} & \textbf{Material} & \textbf{Especificacion} \\
\midrule
OE.5.2.2.4 & Curva PVC pesada (PV-P) 3/4'' & Radio medio 50 mm, angulos 45° y 90°, NTP 399.006 \\
OE.5.2.2.5 & Curva PVC liviana (PV-L) 3/4'' & Radio medio 50 mm, angulos 45° y 90°, NTP 399.006 \\
OE.5.2.2.6 & Curva PVC pesada (PV-P) 1'' & Radio medio 75 mm, angulos 45° y 90°, NTP 399.006 \\
OE.5.2.2.7 & Accesorios de canalizacion (global) & Pegamento, uniones, conectores, abrazaderas, cinta selladora \\
\bottomrule
\end{longtable}

\subsubsection{Requisitos de Materiales}

\begin{longtable}{L{5.0cm} L{2.0cm} L{7.5cm}}
\toprule
\textbf{Material} & \textbf{Cant.} & \textbf{Especificacion tecnica} \\
\midrule
Curva PVC pesada (PV-P) 3/4'' 90° & 35 pza | Para tomacorrientes, radio medio 50 mm \\
Curva PVC pesada (PV-P) 3/4'' 45° & 10 pza | Para tomacorrientes, radio medio 50 mm \\
Curva PVC liviana (PV-L) 3/4'' 90° & 10 pza | Para alumbrado, radio medio 50 mm \\
Curva PVC liviana (PV-L) 3/4'' 45° & 2 pza | Para alumbrado, radio medio 50 mm \\
Curva PVC pesada (PV-P) 1'' 90° & 4 pza | Para alimentador, radio medio 75 mm \\
Curva PVC pesada (PV-P) 1'' 45° & 2 pza | Para alimentador, radio medio 75 mm \\
Pegamento PVC (cemento solvente) & 1/2 gln | Para todo tipo de tuberia PVC electrica \\
Conectores tipo "romex" 3/4'' & 80 pza | Para conexion tuberia-caja \\
Conectores tipo "romex" 1'' & 4 pza | Para conexion tuberia-caja (alimentador) \\
Uniones universales PVC 3/4'' & 30 pza | Para empalme de barras \\
\bottomrule
\end{longtable}

\subsubsection{Metodo de Ejecucion}

\begin{enumerate}
  \item Seleccionar el accesorio del tipo y diametro adecuado a la tuberia.
  \item Aplicar pegamento PVC en ambos elementos a unir.
  \item Ensamblar con presion y giro de 1/4 de vuelta.
  \item Mantener 15 segundos para fraguado inicial.
  \item Para el global de accesorios (OE.5.2.2.7), se considera el pegamento, uniones, conectores y abrazaderas necesarios para completar la instalacion de todas las tuberias del proyecto.
\end{enumerate}

\subsubsection{Control de Calidad}

\begin{itemize}
  \item Verificar que las curvas no presenten reduccion de seccion mayor al 10\%.
  \item Comprobar que el radio de curvatura permita el paso libre de conductores.
  \item Inspeccionar que las uniones esten completamente selladas (sin espacios visibles).
  \item Verificar que no existan rebabas internas en los accesorios.
\end{itemize}

\subsubsection{Metodo de Medicion}

Las curvas se mediran por pieza (pza) instalada. Los accesorios de canalizacion (OE.5.2.2.7) se mediran como un global (glb) por toda la obra.

\subsubsection{Unidad de Medida}

pza (pieza) para curvas; glb (global) para accesorios.

\subsubsection{Forma de Pago}

Pago por pieza de curva instalada o por global de accesorios completado, verificado y aprobado por supervision.
